\begin{topic}{0. Keyword Index by Exam Term}
\begin{tabularx}{\linewidth}{@{}lY@{}}
\toprule
Category & Search these terms / go to sections\\
\midrule
DBMS \& Architecture & DBMS purpose, efficiency, reliability (1); centralised, distributed, cloud, grid, P2P, web (4)\\
Storage \& I/O & HDD, SSD, seek, rotation, cache, buffer (2); RAID, availability, CRC, reliable write/message, WAL, ARIES, checkpoint (24--31)\\
SQL \& Query & SQL, relational algebra, selectivity, cardinality, statistics, enumeration, heuristic, adaptive, rewrite, nested-loop join (5--9)\\
Indexes \& Spatial & index choice, hash, B+ tree, bitmap, R-tree, quadtree, EXPLAIN, fragmentation (10--13)\\
Transactions & ACID, isolation, flat transaction, savepoint, cursor, nested transaction, parent/child, TP monitor, 2PC (14--16, 32)\\
Concurrency \& Locks & dirty read, lost update, conflict, dependency graph, serialisability, wormhole, S/X, 2PL, isolation degree, deadlock, intention lock, update lock, OCC, snapshot (17--23, 33)\\
Distributed Consistency & distributed transaction, replication, available copies, CAP, eventual consistency, BASE, PACELC (32--35)\\
Backup \& Warehouse & remote backup, one-safe/two-safe, full/differential/incremental backup, data warehouse, ETL (36--37)\\
\bottomrule
\end{tabularx}
\end{topic}

\begin{topic}{1. DBMS Purpose, Efficiency and Reliability}
\points
\begin{itemize}
  \item A DBMS stores an integrated collection of data and provides controlled, efficient and reliable access.
  \item \kw{Efficiency}: response time and throughput; affected by storage, I/O, memory, architecture, data model, indexes and query plans.
  \item \kw{Effectiveness}: concurrent users obtain correct results, not merely fast results.
  \item \kw{Security/reliability}: access control, transactions, recovery, replication and fault tolerance preserve data and service.
\end{itemize}
\answer State the required property, identify the workload/failure/concurrency assumption, name the mechanism, then explain its benefit and cost.
\worked ``The system must serve many users after a node failure'' requires availability and fault tolerance; replication helps, but adds coordination and consistency cost.
\end{topic}

\begin{formula}{2. HDD, SSD, Cache and Buffer I/O}
\[
T_{\mathrm{HDD}}=T_{\mathrm{seek}}+T_{\mathrm{rotation}}
+\frac{\mathrm{transfer\ size}}{\mathrm{bandwidth}},
\qquad
T_{\mathrm{SSD}}=\frac{\mathrm{transfer\ size}}{\mathrm{bandwidth}}
\]
\begin{itemize}
  \item Seek moves the head to a track; rotation waits for the sector; transfer moves the bytes.
  \item Large sequential I/O is dominated by transfer time; small random I/O is dominated by seek/rotation.
\end{itemize}
\worked With seek \(12\) ms, rotation \(4\) ms, bandwidth \(4\) MB/s: \(8\) MB takes \(12+4+2000=2016\) ms; \(8\) KB takes about \(18\) ms. The small request spends most of its time positioning the disk.
\[
EA=HC+(1-H)M,\qquad
EA_{\mathrm{buffer}}=H_BB_C+(1-H_B)D
\]
\begin{itemize}
  \item Weight hit and miss paths by probability; compare the final effective time, not hit ratio alone.
  \item Tutorial example: \(0.5C+0.5(100C)=50.5C\), while \(0.9C+0.1(400C)=40.9C\); the larger-cache machine is faster despite slower memory.
\end{itemize}
\trap{Keep units consistent. Convert MB/s to KB/ms before adding millisecond latencies.}
\end{formula}

\begin{topic}{3. NoSQL: Key-Value, Document, Column, Graph}
\begin{tabularx}{\linewidth}{@{}lYY@{}}
\toprule
Model & Strength and suitable workload & Limitation\\
\midrule
Key--value & \(\text{key}\arr\text{value}\); direct lookup, carts, sessions, caching; simple and scalable & Weak joins, relationships and rich predicates\\
Document & JSON-like, nested and varying fields; articles/catalogues and evolving semi-structured records & Cross-document joins and constraints are weaker\\
Column-based & Reads selected columns efficiently; large analytical scans and aggregation & Less natural for frequent whole-record retrieval/update\\
Graph & Vertices and edges; multi-hop traversal, recommendations, friend-of-friend and paths & Extra complexity when relationships are not central\\
\bottomrule
\end{tabularx}
\begin{itemize}
  \item Select from the \textbf{dominant access pattern}, not merely the input format.
  \item NoSQL commonly trades some relational joins, constraints, query flexibility or consistency guarantees for schema flexibility and horizontal scale.
\end{itemize}
\answer Identify the data structure and main operations; choose the model; explain the matching strength; finish with one lost capability.
\trap{NoSQL means ``Not Only SQL'', not ``no SQL''.}
\end{topic}

\begin{topic}{4. Architecture: Centralised, Distributed, Cloud}
\begin{itemize}
  \item \kw{Centralised}: simple administration and consistency; limited scale and one failure domain.
  \item \kw{Distributed}: owned sites share data/work; locality, scale and fault tolerance improve, but global coordination is harder.
  \item \kw{Cloud}: elastic third-party infrastructure and low initial capital cost; privacy, control, vendor dependence and recurring cost matter.
  \item \kw{Grid}: autonomous nodes contribute resources to a task. \kw{P2P}: dynamic decentralised membership. \kw{Web}: globally distributed data with different owners.
\end{itemize}
\worked Startup with one office, little capital and expected rapid growth: cloud is attractive because capacity can scale without building data centres. A large organisation requiring direct ownership/privacy across its own cities may prefer a distributed database.
\answer Compare scale, ownership, privacy, availability, budget and administration. Borderline answers are acceptable when the trade-off is explicit.
\end{topic}

\begin{topic}{5. SQL, Relational Algebra and Query Pipeline}
\[
\begin{aligned}
\text{SQL}&\arr\text{parse/translate}\arr\text{relational algebra}\\
&\arr\text{physical plans}\arr\text{lowest estimated cost}\arr\text{execute}
\end{aligned}
\]
\begin{itemize}
  \item Selection \(\sigma\) chooses rows; projection \(\pi\) chooses columns; product \(\times\) forms all pairs; theta join \(\bowtie_\theta\) applies a condition; natural join \(\bowtie\) uses common attributes.
  \item Example: \texttt{SELECT name FROM E WHERE salary>80000}
  becomes \(\pi_{\mathrm{name}}(\sigma_{\mathrm{salary}>80000}(E))\).
  \item Logical operators say \emph{what}; heap/index scan and nested-loop/hash/merge join say \emph{how}.
  \item Equivalent expressions return the same result but may create very different intermediate sizes.
\end{itemize}
\answer Identify the query result, map clauses to logical operators, propose a useful rewrite, then discuss the physical access path.
\end{topic}

\begin{formula}{6. Selectivity, Cardinality and Statistics}
\[
|\sigma_\theta(R)|\approx RF(\theta)|R|,
\qquad
RF(A=v)\approx\frac{1}{V(A,R)}
\]
\[
RF(\theta_1\land\theta_2)\approx RF(\theta_1)RF(\theta_2)
\quad\text{if independent}
\]
\[
|R\bowtie_{A=B}S|\approx
\frac{|R||S|}{\max(V(A,R),V(B,S))}
\]
\begin{itemize}
  \item \(V(A,R)\): distinct \(A\)-values. Keep tuple count \(n_R\), page count \(b_R\) and distinct count \(V(A,R)\) separate.
  \item Uniformity and independence are assumptions. Skew and correlated predicates can make estimates wrong.
  \item Statistics guide plan selection; sampling and histograms improve estimates.
\end{itemize}
\worked After many insertions/deletions, total rows may be unchanged but distributions/statistics may be stale. The optimiser then misestimates selectivity and chooses a formerly good but now poor plan. Refresh/recompile statistics and inspect the new plan.
\end{formula}

\begin{topic}{7. Query Optimisation: Enumeration, Heuristic, Adaptive}
\begin{itemize}
  \item \kw{Exhaustive enumeration}: cost many equivalent plans and select the cheapest. Best justified for expensive, complex multi-join queries; optimisation overhead/search space can be large.
  \item \kw{Heuristic}: apply rules such as early selection/projection and restrictive operations first. Fast and suitable for simple queries, but not guaranteed globally optimal.
  \item \kw{Adaptive}: execute part of a cost-based plan, observe actual cardinality and choose/revise the remaining plan. Useful when estimates are uncertain.
  \item Real systems often use heuristics to prune the search space before cost-based comparison.
\end{itemize}
\worked Tutorial scenarios: a single-table query with two predicates can use heuristics; a five-table join benefits from cost-based enumeration, possibly adaptive re-estimation. A purely heuristic plan has no cost model to adapt in the same sense.
\answer Compare query complexity, optimisation time, estimation uncertainty and the value of runtime feedback.
\end{topic}

\begin{topic}{8. Query Rewriting, Materialisation and Tuning}
\begin{itemize}
  \item Push selective predicates down: \(\sigma_c(R\bowtie S)\equiv\sigma_c(R)\bowtie S\) when \(c\) uses only \(R\).
  \item Push projections down to reduce tuple width, but retain join, filter and output attributes.
  \item Optimise the complete plan; the cheapest isolated operator may create an expensive intermediate result.
  \item Refresh statistics; add an appropriate index; parameterise/rewrite repeated queries; force a plan only when the optimiser repeatedly fails.
  \item Store derived/pre-joined data when computation is frequent and source changes are infrequent; pay refresh/storage cost.
\end{itemize}
\worked If \(R\) has one million rows but only \(1\%\) satisfy \(R.status='active'\), filtering before the join reduces the candidate rows to about \(10{,}000\).
\end{topic}

\begin{formula}{9. Nested-Loop Join Cost: Tuple vs Page}
Let \(r\) be outer, \(s\) inner, \(n_r\) outer tuples and \(b_r,b_s\) page counts.
\[
\text{Tuple NLJ:}\quad I/O=b_r+n_rb_s,\qquad seeks=b_r+n_r
\]
\[
\text{Page NLJ:}\quad I/O=b_r+b_rb_s,\qquad seeks=2b_r
\]
\begin{itemize}
  \item Test both orientations: making the smaller relation outer often reduces repeated inner scans.
  \item Page NLJ reuses one outer page across an inner scan and is normally much cheaper than tuple NLJ.
  \item Merge join is attractive for ordered inputs; hash join for equality; neither choice should be justified without predicate/order/size assumptions.
\end{itemize}
\answer Define \(n\) versus \(b\), substitute carefully, evaluate both orientations, and state the cheaper plan.
\trap{Under the course convention, tables currently in main memory do not change the optimiser's page-cost calculation. This is a sample-exam rule.}
\end{formula}

\begin{solidtopic}{10. Index Choice: Hash, B+ Tree, Bitmap, R-tree}
\begin{itemize}
  \item \kw{Equality on ID/name}: Hash or B+ tree. Hash gives direct bucket lookup; B+ tree also supports ordered lookup.
  \item \kw{Range/inequality/order}: B+ tree. Ordered leaves support range scans and sorted traversal.
  \item \kw{Low-cardinality filters}: Bitmap. Compact bit vectors support fast AND/OR over values such as status/state.
  \item \kw{Region/distance/nearest neighbour}: R-tree or quadtree. Spatial partitions/boxes prune irrelevant regions.
  \item \kw{Aggregate over every row}: usually no filtering index. Every record still contributes, so a scan may be cheaper.
\end{itemize}
\begin{itemize}
  \item Always give \textbf{index type + search key + predicate + reason}.
  \item Sample exam: income \(<70{,}000\) \arr B+ tree on income; within 5 km \arr R-tree/quadtree on location; average income \arr no filtering index.
  \item A full scan may win for a small table, low-selectivity predicate or query needing most rows.
  \item Clustered B+ tree helps when adjacent key values qualify; unclustered indexes may cause many random page reads.
  \item Composite B+ tree \((A,B,C)\): use leftmost prefixes \(A\), \(A,B\), \(A,B,C\); \(B\) alone normally cannot seek.
  \item \kw{Answer pattern}: state access pattern/selectivity, name the index, explain avoided pages/work, then mention write and space cost.
  \item Equality-only unique lookup favours hash; use B+ tree when range/order/prefix matters. Bitmap suits read-mostly analytics; frequent updates make maintenance unattractive.
\end{itemize}
\worked \texttt{WHERE salary > 0}: likely no filtering index because most rows qualify. \texttt{WHERE id=?}: hash/B+ tree because it is highly selective.
\answer ``Use a B+ tree with search key income because the predicate defines an ordered range; the tree finds the first qualifying key and scans adjacent leaves.''
\end{solidtopic}

\begin{topic}{11. Index Costs, EXPLAIN, Fragmentation, Statistics}
\begin{itemize}
  \item Benefits: fewer page reads for selective lookup; ordered range access; possible join/order/group support; a covering index may avoid base-table reads.
  \item Costs: disk and buffer-pool space; every INSERT adds entries; DELETE removes them; indexed-column UPDATE removes/reinserts entries; writes and administration become slower.
  \item Splits, merges, sparse pages and fragmentation can arise; unused/duplicate indexes provide cost without benefit.
\end{itemize}
\why An indexed UPDATE may still use the index to locate target rows quickly. It then pays base-page writes, logging/locking and index maintenance; if the indexed key changes, the old entry is deleted and a new entry inserted.
\kw{Diagnosing slowdown after insert/delete}
\begin{itemize}
  \item \textbf{Stale statistics}: distributions changed, so cardinality estimates and the selected plan are wrong. Refresh statistics and compare estimates with actual rows.
  \item \textbf{Fragmentation/sparse pages}: splits and deletions increased pages or random I/O. Measure it, then reorganise/rebuild when material.
  \item The unchanged row count does not imply unchanged distribution, page layout, clustering or cache behaviour.
\end{itemize}
\texttt{CREATE INDEX idx ON R(A); \quad DROP INDEX idx;}

\texttt{SELECT DISTINCT title FROM titles;}

\texttt{SELECT COUNT(title) FROM titles WHERE title='Engineer';}
\begin{itemize}
  \item \texttt{EXPLAIN}: estimated access method, index, rows, filtering and join order.
  \item \texttt{EXPLAIN ANALYZE}: executes and adds actual rows, time and loops. Large estimate error suggests stale statistics/skew/correlation.
  \item \texttt{Using where}: filtering occurs. \texttt{Using index}: possibly covering; it does not prove a small scan.
\end{itemize}
\answer Compare the same query/data/cache conditions. Identify the changed operator, first large estimated-versus-actual row divergence, avoided pages/loops/sort, and measured runtime/buffer effect. Do not claim improvement from the index name alone.
\end{topic}

\begin{formula}{12. B+ Tree: Height, Search, Insert/Delete}
\[
h_{\max}=\left\lceil\log_{\lceil n/2\rceil}K\right\rceil
\]
\begin{itemize}
  \item \(n\): maximum children; \(K\): search-key values. High fan-out keeps the tree shallow; every root-to-leaf path has equal depth.
  \item Internal nodes route; linked leaves contain ordered keys/record pointers.
  \item Equality descends to one leaf. Range finds the first key then follows leaf links.
  \item Insert: leaf insertion, split on overflow, propagate separator; root split increases height. Delete: borrow/redistribute, otherwise merge and propagate.
\end{itemize}
\worked \(K=10{,}000{,}000\): \(n=4\Rightarrow\lceil\log_2K\rceil=24\); \(n=100\Rightarrow\lceil\log_{50}K\rceil=5\). For key ``Crick'', compare at each separator, descend left of Mozart and Einstein, find Crick in the leaf, then follow its record pointer.
\end{formula}

\begin{topic}{13. Hash, Bitmap, Quadtree and R-tree}
\kw{Hash}
\begin{itemize}
  \item Hash key to bucket; excellent equality performance when distribution is uniform; no useful key order for ranges.
\end{itemize}
\kw{Bitmap}
\begin{itemize}
  \item One bitmap per value, one bit per record. For states \([VIC,NSW,WA,VIC,NSW]\):
  \(VIC=10010,\ NSW=01001,\ WA=00100\).
  \item Combine predicates with bitwise operations, e.g. \(VIC\lor NSW=11011\).
\end{itemize}
\kw{Quadtree/R-tree}
\begin{itemize}
  \item Quadtree recursively divides space. R-tree groups objects in minimum bounding rectangles.
  \item Best-first nearest neighbour: prioritise the node/object with smallest minimum distance; stop only when the nearest queued object cannot be beaten.
  \item Overlapping R-tree boxes may require multiple branches. Tutorial example: point \(i\) lies in BB1 and BB3, so both zero-distance branches must be expanded; object G is then nearest.
\end{itemize}
\end{topic}

\begin{topic}{14. Transactions, ACID and Isolation}
\begin{itemize}
  \item \kw{Atomicity}: all actions occur or none. \kw{Consistency}: constraints hold before/after a correct transaction.
  \item \kw{Isolation}: concurrent execution has the effect of some serial execution, without requiring literal one-at-a-time execution.
  \item \kw{Durability}: committed effects survive failure.
  \item Serial execution is correct but wastes parallel CPU/I/O and lets long transactions block everyone; isolation seeks serial-equivalent results with concurrency.
\end{itemize}
\kw{Distinguish the properties}: atomicity asks whether all effects survive together; isolation asks what concurrent/intermediate effects others may observe; durability asks whether a committed result survives a crash.
\worked A transfer debits \(A\) and credits \(B\). Atomicity prevents debit without credit; consistency preserves the total-balance invariant; isolation hides the half-completed transfer; durability preserves both updates after commit.
\answer Define the property, name the anomaly/failure it prevents, then explain how concurrency/recovery mechanisms enforce it.
\end{topic}

\begin{topic}{15. Flat Transactions, Savepoints and Cursor}
\begin{itemize}
  \item \kw{Flat}: one unit; failure can discard all work. Poor for a million-row bonus update because one late error wastes a long transaction.
  \item \kw{Savepoints}: rollback to an intermediate state; earlier work is retained but becomes durable only at final commit.
  \item \kw{Chained}: linked committed segments. \kw{Nested}: parent invokes subtransactions.
\end{itemize}
\texttt{BEGIN WORK; \quad COMMIT WORK; \quad ROLLBACK WORK;}
\begin{itemize}
  \item Tutorial trace rule: restore the exact variable/database state recorded at the named savepoint, then continue from there.
  \item A \kw{singleton SELECT} returns one row directly; a \kw{cursor} points into a multi-row result and fetches rows one at a time.
\end{itemize}
\worked If a rollback returns to savepoint 2 where \(AccBalance1=100\), later changes setting it to \(0\) disappear; a second rollback to savepoint 2 restores \(100\) again.
\end{topic}

\begin{topic}{16. Nested Transactions, Parent/Child and TP Monitor}
\kw{Nested rules}
\begin{itemize}
  \item Child commit is provisional: durable only when every ancestor commits.
  \item If a subtransaction rolls back, all its descendants roll back.
  \item Child changes become visible to its parent after child commit; parent objects are visible to children.
\end{itemize}
\worked If A, B and C commit but PARENT aborts, all roll back. If A, B and PARENT commit while C aborts, A/B/PARENT are durable and C is rolled back.
\kw{TP monitor}
\begin{itemize}
  \item Integrates clients, services and heterogeneous DBMSs; coordinates communication, context, restart and global transaction work.
  \item One process for every service is a bottleneck and concentrated failure point; distribution improves performance and fault containment.
\end{itemize}
\end{topic}

\begin{topic}{17. Dirty Read, Lost Update and Conflict Types}
\begin{itemize}
  \item \kw{Dirty read}: read an uncommitted value that may be rolled back.
  \item \kw{Unrepeatable read}: two reads of the same object see an intervening committed write.
  \item \kw{Lost update}: one write overwrites a result computed from an earlier value by another transaction.
\end{itemize}
\begin{tabularx}{\linewidth}{@{}lY@{}}
\toprule
Order on object \(O\) & Dependency\\
\midrule
\(W_i(O)\arr W_j(O)\) & write--write\\
\(W_i(O)\arr R_j(O)\) & write--read\\
\(R_i(O)\arr W_j(O)\) & read--write\\
\(R_i(O)\arr R_j(O)\) & no conflict dependency\\
\bottomrule
\end{tabularx}
\worked With \(B=50\), \(T_1:\times2,\times2\) and \(T_2:+10\): serial outcomes are \(210\) or \(240\). Interleaving after the first multiply can produce \(220\), which matches neither serial order and violates isolation.
\end{topic}

\begin{formula}{18. Dependency Graph, Serialisability, Wormhole}
If \(T_i\) performs a conflicting action on \(O\) before \(T_j\):
\[
(T_i,O,T_j)\in DEP(H)
\]
\begin{enumerate}
  \item Group actions by object and preserve time order.
  \item Ignore read--read; add directed edges for W--W, W--R and R--W.
  \item A directed cycle means no equivalent serial history. A DAG topological order gives a serial order; more than one topological order may be valid.
\end{enumerate}
\[
Before(T)=\{U\mid U\leadsto T\},\qquad After(T)=\{U\mid T\leadsto U\}
\]
\[
Before(T)\cap After(T)\ne\varnothing
\iff T\text{ is in a cycle/wormhole}
\]
\worked Tutorial cycle \(T_1\arr T_3\arr T_2\arr T_1\) is not serialisable. Reversing the \(O_1\) edge gives \(T_3\arr T_2\arr T_1\), a valid serial order (independent \(T_5\) may be placed consistently).
\worked If \(W_1(A)\) precedes \(R_2(A)\), add \(T_1\arr T_2\). If \(W_2(B)\) precedes \(R_1(B)\), add \(T_2\arr T_1\). The two-edge wormhole proves non-serialisability.
\answer List edges with their object/conflict, draw/check the graph, state cycle or topological order, then conclude isolated/not isolated.
\end{formula}

\begin{topic}{19. S/X Locks, Well-Formedness and 2PL}
\begin{tabularx}{\linewidth}{@{}lcc@{}}
\toprule
Current \(\backslash\) Request & S & X\\
\midrule
Free & \yes & \yes\\
S & \yes & \no\\
X & \no & \no\\
\bottomrule
\end{tabularx}
\begin{itemize}
  \item READ needs S or stronger; WRITE needs X.
  \item \kw{Well formed}: every READ/WRITE/UNLOCK is covered by an earlier appropriate lock.
  \item \kw{Two phase}: growing phase acquires locks; after the first unlock, shrinking begins and no new lock may be acquired.
  \item If all transactions are well formed and two phase, every legal lock history is isolated. 2PL ensures serialisability but may deadlock.
\end{itemize}
\kw{Three separate checks}: legal means currently granted locks are compatible; well formed means operations hold the required lock; 2PL means no acquisition/upgrade after the first release/downgrade. A history can pass one and fail another.
\why Each transaction has a lock point at its final acquisition. Ordering transactions by lock point is consistent with their conflicts, preventing a dependency cycle.
\kw{Variants}: strict 2PL holds X locks to commit/abort and prevents cascading aborts; rigorous 2PL holds both S and X locks. Neither eliminates deadlock.
\worked In a numbered trace, mark the first UNLOCK; every later LOCK line violates 2PL. Per-operation lock/unlock alone is insufficient because releasing between related operations can still produce a non-serial result.
\answer ``The history is legal, but \(T_2\) requests \(X(B)\) after releasing \(S(A)\), so \(T_2\) violates 2PL and the 2PL serialisability theorem cannot be applied.''
\end{topic}

\begin{topic}{20. Isolation Degrees 0--3}
\begin{tabularx}{\linewidth}{@{}cYY@{}}
\toprule
Degree & Guarantee & Required locking\\
\midrule
0 & Weak conflict protection & Writes well formed\\
1 & No lost updates & X locks well formed and 2PL; reads may be unlocked\\
2 & No lost updates/dirty reads; unrepeatable reads possible & Reads/writes well formed; X locks 2PL; S may release after read\\
3 & Repeatable reads/full conflict protection & Complete well-formed 2PL for S and X\\
\bottomrule
\end{tabularx}
\[
\begin{array}{c|c}
D1&W\arr W\\
D2&W\arr W,\ W\arr R\\
D3&W\arr W,\ W\arr R,\ R\arr W
\end{array}
\]
\answer Find the weakest operation: an unlocked READ prevents Degree 2/3; an X-lock sequence that is not two phase prevents Degree 1+.
\end{topic}

\begin{topic}{21. Deadlock, Livelock, Spin Locks, Semaphores}
\begin{itemize}
  \item Deadlock is a cycle in a wait-for/resource-dependency graph; each transaction holds something the next needs.
  \item Prevent with global resource order or preallocation; detect cycles and roll back a victim; timeout is practical for distributed deadlocks.
  \item Immediate repeated rollback can cause livelock. Risk rises with concurrency, locks per transaction, hold time, hotspots and write conflicts.
  \item OS locks sleep but pay context-switch cost; spin locks use atomic test-and-set/CAS and suit short low-contention waits; semaphores queue waiters.
  \item FIFO handoff can create a convoy; waking all waiters reduces serial handoff but causes competition.
\end{itemize}
\worked Tutorial approximation with \(n\) other transactions, \(r\) exclusive locks each and \(R\) records gives lifetime wait probability \(\approx nr^2/(2R)\) and pair-deadlock probability growing roughly with \(n^2r^4/(4R^2)\). Use this only under the stated simplifying assumptions.
\end{topic}

\begin{formula}{22. Granular Locks: IS, IX, SIX, U, X}
\scriptsize
\begin{tabularx}{\linewidth}{@{}l|cccccc@{}}
Request \(\backslash\) Current & IS & IX & S & SIX & U & X\\
\hline
IS  & \yes&\yes&\yes&\yes&\yes&\no\\
IX  & \yes&\yes&\no &\no &\no &\no\\
S   & \yes&\no &\yes&\no &\no &\no\\
SIX & \yes&\no &\no &\no &\no &\no\\
U   & \yes&\no &\yes&\no &\no &\no\\
X   & \no &\no &\no &\no &\no &\no
\end{tabularx}
\normalsize
\begin{itemize}
  \item Rows are requests; columns are current modes. The matrix is directional for U/S in this course.
  \item IS intends finer S; IX intends finer S/X; SIX is S here plus IX below; U reads with intent to update.
  \item Acquire root-to-leaf, release leaf-to-root. S/IS on a non-root needs a suitable parent; X/U/SIX/IX needs IX or stronger on every ancestor.
\end{itemize}
\kw{Request procedure}
\begin{enumerate}
  \item Process stated arrival order and record only locks already \textbf{granted}.
  \item Verify the required intention mode on every ancestor.
  \item Compare the request with every granted target lock using the matrix.
  \item Grant only if all checks pass; otherwise wait. A waiting request is not a current lock.
  \item State queue policy: strict FIFO may block a later compatible request behind an earlier incompatible one.
\end{enumerate}
\answer For every request state ancestor status, current granted modes, compatibility result, grant/wait decision and resulting lock table. A compatible leaf lock cannot override a root conflict.
\worked Sample pattern: T1's IS and T2's SIX can coexist at root, so T1/T2 run; later T3's S conflicts with existing SIX and waits. If T3 arrives before T2, T3's S can instead make T2 wait. Arrival order changes the runnable set.
\end{formula}

\begin{topic}{23. Update Locks, Escalation, OCC, Snapshot}
\begin{itemize}
  \item U lock admits one intended updater while allowing compatible earlier readers under the course matrix; it reduces S-to-X conversion deadlock.
  \item Lock escalation replaces many fine locks with a coarse S/X lock: less management overhead, larger conflict range and less concurrency.
  \item \kw{Pessimistic/2PL}: conflicts expected; prevent them before access; cost is locks, waits and deadlocks.
  \item \kw{Optimistic}: conflicts rare; read/compute freely, validate near commit; cost is wasted work/retry when conflicts are common.
  \item \kw{Snapshot isolation}: versioned consistent reads and first-writer-commits; high read concurrency but write skew can violate cross-row constraints.
  \item \kw{Timestamp ordering}: globally consistent order; distributed form carries a coordinator-assigned timestamp to every server.
\end{itemize}
\worked Two transactions read \(A=B=100\), then independently set \(A=-100\) and \(B=-100\). No direct write conflict occurs, yet \(A+B\ge0\) fails: snapshot isolation is not automatically serialisable.
\end{topic}

\begin{formula}{24. Probability, Availability, RAID 0/1/5/6}
\[
P(A\cap B)=P(A)P(B)
\]
\[
P(A\cup B)\approx P(A)+P(B)
\quad\text{for small independent probabilities}
\]
\[
Availability=\frac{MTTF}{MTTF+MTTR},\qquad
P(\text{unavailable})\approx\frac{MTTR}{MTTF}
\]
If failure requires \(k\) of \(n\) disks simultaneously:
\[
P_{\mathrm{fail}}\approx\binom{n}{k}p^k,\qquad
MTTF_{\mathrm{system}}\approx\frac{M^k}{\binom{n}{k}}
\]
\scriptsize
\begin{tabularx}{\linewidth}{@{}cYccc@{}}
\toprule
RAID & Organisation & Tolerates & Capacity & Course example\\
\midrule
0 & Block striping & 0 & \(nS\) & 2 disks: \(M/2\)\\
1 & Mirroring & \(n-1\) & \(S\) & 2: \(M^2\), 3: \(M^3\)\\
3/4/5 & Single parity & 1 & \((n-1)S\) & 3 disks: \(M^2/3\)\\
6 & Dual parity & 2 & \((n-2)S\) & 5 disks: \(M^3/10\)\\
\bottomrule
\end{tabularx}
\normalsize
\begin{itemize}
  \item RAID 0: throughput/capacity, no tolerance. RAID 1: best simple redundancy, lowest utilisation. RAID 4: parity-disk write bottleneck. RAID 5 distributes parity. RAID 6 pays extra parity/write cost for two failures.
  \item Tutorial utilisation: RAID 0 \(100\%\), RAID 1 with two disks \(50\%\), RAID 3 with three disks \(2/3\).
  \item Derive reliability by finding the minimum simultaneous failures that lose data, counting such combinations, multiplying by \(p^k\), then inverting the failure rate under the course approximation.
\end{itemize}
\worked Six 2-TB disks provide 12 TB raw. RAID 1 provides 6 TB usable; RAID 5 provides \(5\times2=10\) TB. Capacity and reliability are separate calculations.
\worked RAID 0 loses data after any one of \(n\) disks fails, so \(P_{\rm loss}\approx np\). A mirrored pair requires both failures, \(p^2\). Single parity tolerates one failure but loses data for any pair, \(\binom n2p^2\).
\trap{Tutorial text has a typo for RAID 0; doubled failure probability means MTTF halves, so use \(M/2\). State the independence/small-\(p\) assumption.}
\end{formula}

\begin{topic}{25. Fault Tolerance: Failvote, Failfast, CRC}
\kw{Failvote vs failfast}
\begin{itemize}
  \item Failvote requires a majority of the original modules: stronger confidence, lower availability as failures accumulate.
  \item Failfast excludes unavailable modules and takes a majority among working modules: higher availability, but one remaining module cannot be checked.
\end{itemize}
\kw{Reliable block write}
\begin{itemize}
  \item Duplex write stores two sequential copies with version numbers; choose the newest consistent copy.
  \item Logged write stores the second representation in a log and is efficient for small changes.
  \item CRC detects storage/transmission corruption and burst errors.
\end{itemize}
\answer Separate \emph{correctness of an accepted result} from \emph{availability of the service}.
\end{topic}

\begin{formula}{26. Reliable Message Passing: Out, Ack, In}
\[
Out=\#sent,\qquad Ack=\#acknowledged,\qquad In=\#received
\]
When B sends message \(k\) to A:
\begin{enumerate}
  \item B persistently records \(Out_k\) and sets \(Out=k\).
  \item B sends; A persistently records \(In_k\) and sets \(In=k\).
  \item A acknowledges \(k\); B persistently sets \(Ack=k\).
\end{enumerate}
\worked Initial A \((Out=3,Ack=3,In=6)\), B \((Out=6,Ack=6,In=3)\). B sends message 7 and receives the acknowledgement:
\[
A:(3,3,7),\qquad B:(7,7,3)
\]
\answer First identify sender and receiver. Only sender changes Out/Ack; only receiver changes In. Preserve stable-storage update order so retransmission/recovery is possible.
\end{formula}

\begin{topic}{27. Buffer Pool, STEAL/NO-FORCE, WAL}
\begin{itemize}
  \item \texttt{fix(page)} loads/pins; a fixed page cannot be evicted. \texttt{unfix} releases one user's pin. Dirty means memory is newer than disk.
  \item STEAL allows uncommitted dirty pages to reach disk \arr recovery needs UNDO/old values.
  \item NO-FORCE allows commit before every dirty page reaches disk \arr recovery needs REDO/new values.
  \item ARIES supports efficient STEAL + NO-FORCE.
\end{itemize}
\kw{Write-ahead logging}
\begin{enumerate}
  \item Before a data page is written, its update log must be durable:
  \(\boxed{pageLSN\le flushedLSN}\).
  \item Before commit returns, the log through the commit record must be durable.
\end{enumerate}
\why Rule 1 ensures an on-disk uncommitted update can be undone; Rule 2 ensures a committed update absent from data pages can be redone.
\end{topic}

\begin{topic}{28. Log Records: LSN, pageLSN, CLR}
\begin{tabularx}{\linewidth}{@{}lY@{}}
\toprule
Field & Meaning\\
\midrule
LSN & Increasing log-record identifier\\
pageLSN & Latest update/CLR applied to a page\\
flushedLSN & Largest durable LSN\\
prevLSN/lastLSN & Previous/latest log for a transaction\\
recLSN & Earliest potentially unpropagated update that dirtied a page\\
undoNextLSN & CLR pointer to the next action requiring undo\\
\bottomrule
\end{tabularx}
\begin{itemize}
  \item UPDATE: log XID/page/old/new/prevLSN; update lastLSN/pageLSN; first dirtying adds DPT entry with recLSN.
  \item COMMIT: log commit, flush through commit LSN, return, later write END.
  \item ABORT: log abort, follow prevLSN backwards, undo updates and write CLRs, then END.
  \item CLR is redo-only: it records an undo already performed and is never undone again.
\end{itemize}
\end{topic}

\begin{warningbox}{29. ARIES: Analysis, Redo, Undo}
\[
\boxed{\text{Analysis}\arr\text{Redo (repeat history)}\arr\text{Undo losers}}
\]
\kw{Analysis}
\begin{itemize}
  \item Start from checkpoint via master record; initialise transaction table and DPT from END\_CHECKPOINT.
  \item Scan forward: update status/lastLSN; first UPDATE/CLR for a missing page adds DPT recLSN. END removes a transaction.
  \item Committed/ended work is winner work; active/aborting transactions remaining are losers.
\end{itemize}
\kw{Redo}
\begin{itemize}
  \item Start at smallest DPT recLSN. For each UPDATE/CLR, skip if page absent from DPT, \(recLSN>LSN\), or disk \(pageLSN\ge LSN\); otherwise reapply and set pageLSN.
  \item Redo repeats crash-time history, including loser updates/CLRs; it is not ``redo committed only''.
\end{itemize}
\smallskip
\kw{Undo}
\begin{itemize}
  \item Initialise \(ToUndo\) with loser lastLSNs; repeatedly take the largest.
  \item Ordinary update: restore old value, write CLR, follow prevLSN. CLR: follow undoNextLSN. End of chain: write END.
\end{itemize}
\why ARIES first restores the exact physical state at crash, then undoes losers once in a controlled order. A CLR records completed undo and is redo-only, so a second crash does not repeat that undo incorrectly.
\kw{Interpret the structures}: the transaction table tracks status/log-chain end; DPT recLSN is the earliest update that may be missing from disk for a dirty page; pageLSN shows the latest update already present on that disk page.
\end{warningbox}

\begin{topic}{30. ARIES Worked Pattern: DPT, TT, ToUndo}
\worked Tutorial log:
\[
\begin{array}{rll}
10&T1:&P1:Y\arr Z\\
15&T2:&P3:U\arr V\\
25&CP:&TT=[T1:10,T2:15],DPT=[P1:10,P3:15]\\
30&T1:&P2:W\arr X\\
35&T1:&COMMIT\\
40&T2:&P1:Z\arr T\\
45&T2:&ABORT\\
50&T2:&CLR\ P1=Z,\ undoNext=15
\end{array}
\]
\begin{itemize}
  \item Analysis reconstructs \(P1:10,P3:15,P2:30\); T1 is committing/winner, T2 aborting/loser with lastLSN 50.
  \item Redo starts at 10 and conditionally reapplies 10,15,30,40 and CLR 50 according to DPT/pageLSN.
  \item Undo starts \(ToUndo=\{50\}\); CLR 50 points to 15; undo LSN 15 by restoring \(P3=U\), write a CLR, then END.
\end{itemize}
\answer Present the final Analysis transaction table/DPT; a Redo row for each candidate showing DPT, recLSN and pageLSN tests; and a descending ToUndo trace showing each CLR, undoNextLSN and next queue state.
\end{topic}

\begin{topic}{31. Checkpoint, Shadow Paging, Recovery Cost}
\kw{Checkpoint trade-off}
\begin{itemize}
  \item Frequent checkpoints reduce the log/history considered after a crash and shorten recovery.
  \item They add normal-time log entries, disk I/O and transaction overhead. Choose based on recovery urgency versus throughput.
\end{itemize}
\kw{Shadow paging}
\begin{itemize}
  \item Keep immutable shadow and current page tables; update pages by copy-on-write.
  \item Commit flushes modified pages/current table, then atomically redirects the fixed shadow pointer.
  \item Advantage: no update log and simple crash recovery. Costs: page-table copying, garbage collection, fragmentation, high commit flushing and weak concurrency support.
\end{itemize}
\end{topic}

\begin{topic}{32. Distributed Transactions and 2PC}
\begin{itemize}
  \item Atomic commit and global serial order are different: 2PC gives all-commit/all-abort; locking/timestamps/optimistic validation give consistent conflict order.
  \item \kw{Phase 1}: coordinator sends PREPARE; participants force prepare state and vote YES/NO.
  \item \kw{Phase 2}: all YES \arr record/send GLOBAL COMMIT; any NO/timeout \arr GLOBAL ABORT; participants record decision, act locally and release locks.
  \item A YES vote promises that the participant has durable prepare state and can commit later; it therefore cannot unilaterally abort while uncertain.
\end{itemize}
\worked
\begin{itemize}
  \item P1/P2 YES, P3 NO \arr abort all.
  \item Participant crashes before voting \arr timeout and abort.
  \item Coordinator crashes after all YES \arr prepared participants cannot decide independently; request the decision and may block until recovery.
\end{itemize}
\why Blocking preserves atomicity: committing locally could conflict with another participant's abort. The cost is retained locks/resources and reduced availability while the durable global decision is unavailable.
\answer State the durable state already reached, who can still decide, the timeout/recovery action, and why locks/blocking remain.
\end{topic}

\begin{topic}{33. Distributed Concurrency, Replication, Copies}
\begin{itemize}
  \item Locking must preserve the same T-before-U conflict order at every relevant server.
  \item Timestamp ordering uses a globally unique coordinator timestamp carried with each object access.
  \item Optimistic distributed control validates at each server during 2PC phase 1 using a globally consistent validation order.
  \item Available Copies reads one available replica and writes all available replicas, but changing availability sets can hide conflicts and break one-copy serialisability.
  \item Final conservative rule: before commit, verify that the contacted failed/available set did not change; otherwise abort, even if the particular execution might have been harmless.
\end{itemize}
\kw{Why local correctness is insufficient}: each server may have a serialisable local schedule but choose a different order for the same two global transactions. If site \(A\) orders \(T\arr U\) and site \(B\) orders \(U\arr T\), the combined dependency graph has a cycle.
\worked Under Available Copies, \(T_1\) may read/write replica \(A\) while \(B\) is considered failed. If membership changes and \(T_2\) later uses stale \(B\) without learning \(T_1\)'s update, both transactions can appear locally valid but violate one-copy behaviour.
\answer Explain the desired one-logical-copy behaviour, show how sites could miss one another, then state the validation/abort rule.
\end{topic}

\begin{topic}{34. CAP: Consistency, Availability, Partition}
\begin{itemize}
  \item \kw{Consistency}: one current logical value. \kw{Availability}: every request to a non-failing node receives a response. \kw{Partition tolerance}: service continues despite network separation.
  \item During a partition, strict consistency and availability cannot both be guaranteed.
  \item Strong consistency requires waiting, coordination or rejecting operations; this raises latency and lowers availability.
  \item Eventual consistency allows temporary stale reads but replicas converge if updates stop.
\end{itemize}
\worked Sample exam: many geographically distributed servers, unreliable links and millions of users. Partitions are realistic; serving users with low latency/high availability is important, so eventual consistency is suitable if temporary staleness is acceptable.
\kw{Answer-ready detail}
\begin{itemize}
  \item CP rejects/delays some partitioned requests to avoid conflicting or stale results.
  \item AP continues serving but permits temporary divergence; versions, timestamps, quorums, CRDTs or application rules reconcile replicas later.
  \item State the invariant: banking/payment often cannot accept divergence; feeds, caches or carts may tolerate it.
\end{itemize}
\worked Two partitioned cart replicas accept different additions. Availability is preserved and carts temporarily differ; set-union merge can converge later if duplicate/removal semantics are defined.
\answer Identify partition risk; explain the C/A conflict; connect the choice to user-visible behaviour; state convergence/conflict resolution and its cost.
\end{topic}

\begin{formula}{35. Eventual, Causal, BASE and PACELC}
\begin{itemize}
  \item Strong: every read sees latest. Eventual: convergence later. Causal: preserve cause-before-effect.
  \item Read-your-writes: see own latest update. Session: guarantee within one session. Monotonic read: never go backwards. Monotonic write: one process's writes keep issue order.
\end{itemize}
\[
\text{BASE}=\text{Basically Available + Soft State + Eventual Consistency}
\]
\[
\boxed{\text{If partition: A or C;\quad Else: Latency or C}}
\]
\begin{itemize}
  \item PA/EL: available under partition, low latency normally. PC/EC: consistency in both cases.
  \item PA/EC: available under partition, consistency normally. PC/EL: consistent under partition, low latency normally.
  \item CAP describes the forced choice during partition; PACELC additionally describes latency-versus-consistency during normal operation.
  \item NoSQL/partitioning link: scalable NoSQL systems often shard data and use BASE-style relaxed consistency to gain availability, latency and partition tolerance.
\end{itemize}
\worked Shopping cart/product display may favour availability; checkout/payment needs stronger consistency. Airline booking can tolerate more staleness with many seats than near sell-out. Apply guarantees per business function, not universally.
\answer Name the weakest model preserving the application invariant; state the anomaly it permits, the availability/latency benefit, and why the anomaly is acceptable or repairable.
\end{formula}

\begin{topic}{36. Remote Backup: One-Safe, Two-Safe}
\begin{itemize}
  \item Heartbeats/multiple links distinguish primary failure from link failure; bad detection can create split brain.
  \item Backup takeover redoes completed work and removes incomplete work. Periodic redo/checkpointing reduces takeover delay; a hot spare continually applies redo.
  \item \kw{One-safe}: primary log only; low latency, possible acknowledged-data loss before replication.
  \item \kw{Two-very-safe}: primary and backup durable before commit; strongest durability, but site/link failure blocks commit.
  \item \kw{Two-safe}: two-very-safe normally, primary-only during backup outage; better availability but weaker degraded-mode durability.
\end{itemize}
\kw{Failure-analysis pattern}
\begin{itemize}
  \item Identify the last log record known durable at primary and backup.
  \item Decide whether the client already received commit acknowledgement.
  \item On takeover, redo durable committed work and undo/remove incomplete work.
  \item Prevent both sites accepting writes simultaneously using a quorum, lease, fencing token or authoritative membership decision.
\end{itemize}
\worked If one-safe acknowledges after only the primary log is forced, immediate primary loss before shipping can lose an acknowledged transaction. Two-very-safe prevents that loss but cannot acknowledge while the backup/link is unavailable.
\answer Compare durability, availability, latency and behaviour during link/site failure.
\end{topic}

\begin{topic}{37. Backup: Full/Differential, Warehouse, ETL}
\kw{Backup}
\begin{itemize}
  \item Full versus partial depends on change scope; differential stores changes since the last full backup.
  \item Incremental stores changes since the previous backup: smallest daily backup, but restore may need the full backup plus a long chain. Differential grows over time but restore needs only the full plus latest differential.
  \item Higher frequency reduces possible data loss (RPO) but adds I/O/storage; restore design determines recovery time (RTO).
  \item Historical copies enable point-in-time recovery.
\end{itemize}
\kw{Data warehouse}
\begin{itemize}
  \item Integrates historical data under a unified schema for analytics and moves decision support away from OLTP.
  \item Source-driven pushes updates; destination-driven pulls. Exact 2PC synchronisation is usually too costly, so some staleness is accepted.
  \item Requires schema integration, cleansing, deduplication and time dimensions; stored aggregates can replace detail when totals answer the workload.
  \item OLTP favours current detailed records, short updates and concurrency; a warehouse favours historical snapshots, scans, aggregation and stable analytical schemas.
\end{itemize}
\worked A nightly differential plus weekly full backup limits the restore set to two backups, while log archiving between them supports point-in-time recovery. The schedule must still satisfy the required RPO and tested restore time.
\answer Name the objective, choose the mechanism/frequency, and state freshness, recovery, storage and operational trade-offs. For warehousing, also explain why asynchronous ETL and accepted staleness protect OLTP throughput.
\end{topic}
